\section{Methodology}

\subsection{Experimental Setup and Hardware Determinism}
To isolate the physical characteristics of the sensor and prevent systematic measurement errors, 
the data acquisition (DAQ) architecture was strictly engineered for hardware-level determinism. 
An ESP32 microcontroller was utilized exclusively as a localized DAQ unit. 
To eliminate thermal coupling and extraneous Joule heating from the System-on-Chip (SoC) that could compromise microclimate readings, 
all radio frequency (RF) and wireless capabilities were explicitly disabled at the firmware level. 

The primary sensing element, an AHT10 capacitive polymer and silicon sensor, 
was queried via hardware-timer interrupts executing on Core 0. This implementation achieved a highly deterministic sampling interval of $\Delta t = 2.051$\,s, 
with an ultra-low temporal jitter of $<0.05$\,ms. The 2-second nominal polling rate was not arbitrary; 
it was strategically enforced based on hardware design constraints \cite{aht10_datasheet} to restrict internal self-heating to below $0.1^\circ$C during continuous operation, 
a necessity given that relative humidity measurements are directly dependent on the internal thermal state of the sensor.

\subsection{Protocol I: Steady-State Noise Characterization}
To evaluate the intrinsic hardware noise floor ($\sigma$), 
the sensor was mechanically isolated within a thermally insulated enclosure, 
successfully suppressing external convective currents and ambient thermal fluctuations. 
The steady-state acquisition spanned approximately two hours. 

An initial 30-minute segment of the dataset was explicitly trimmed. 
Direct inspection verified that this period represented a physical cooling transient (decaying from $30.99^\circ$C to $30.56^\circ$C), 
resulting from the dissipation of residual heat introduced during sensor handling and installation, rather than active chamber heating. 
Subsequent to this trimming, linear detrending was applied to the near-steady-state data to mathematically decouple the macroscopic environmental drift from the microscopic analog noise inherent to the hardware.

\subsection{Protocol II: Transient Response and Calculus-based Cost Function}
A secondary protocol evaluated the physical phase lag introduced during transient phenomena. 
Following a 3-minute baseline stabilization phase, an isolated thermal step input (warm water injection, $40^\circ\text{C}$--$50^\circ\text{C}$) was introduced into the closed system, 
inducing an instantaneous physical perturbation. 

To condition the analog noise without severely degrading the transient response, 
an Exponential Moving Average (EMA) filter was implemented based on digital signal processing principles \cite{pieterp_ema}:
\begin{equation}
y_t = \alpha \cdot x_t + (1 - \alpha) \cdot y_{t-1} \label{eq:ema02}
\end{equation}
where $x_t$ represents the raw transducer measurement and $y_t$ denotes the filtered signal. 
The theoretical noise power ratio for zero-mean white noise passing through this filter is expressed as:
\begin{equation}
\frac{\sigma_y^2}{\sigma_x^2} = \frac{\alpha}{2-\alpha} \label{eq:noise_power_ratio}
\end{equation}
Consequently, the relative reduction in noise standard deviation ($\text{Red}_{\sigma}$) is modeled as:
\begin{equation}
\text{Red}_{\sigma} = \left(1 - \sqrt{\frac{\alpha}{2-\alpha}}\right) \times 100\% \label{eq:red_sigma}
\end{equation}

The optimal smoothing factor $\alpha$ was derived through numerical optimization of a calculus-based Cost Function $J(\alpha)$. 
This function symmetricized the Min-Max normalized penalty of empirical noise standard deviation ($\hat{C}_{\text{noise}}(\alpha)$) evaluated from Protocol I against the normalized transient phase lag ($\hat{C}_{\text{lag}}(\alpha)$) captured in Protocol II:
\begin{equation}
J(\alpha) = W_1 \cdot \hat{C}_{\text{noise}}(\alpha) + W_2 \cdot \hat{C}_{\text{lag}}(\alpha) \label{eq:cost_func}
\end{equation}
where $W_1 = 0.5$ and $W_2 = 0.5$. The empirical noise cost $\hat{C}_{\text{noise}}(\alpha)$ is defined as the feature-scaled standard deviation of the filtered detrended signal $y_t(\alpha)$:
\begin{equation}
\hat{C}_{\text{noise}}(\alpha) = \frac{\sigma_y(\alpha) - \min_{\alpha}(\sigma_y)}{\max_{\alpha}(\sigma_y) - \min_{\alpha}(\sigma_y)} \label{eq:norm_noise}
\end{equation}
where $\sigma_y(\alpha) = \text{std}(y_t(\alpha))$. Theoretically, for pure zero-mean white noise, the underlying noise power ratio follows $\frac{\sigma_y^2}{\sigma_x^2} = \frac{\alpha}{2-\alpha}$, yielding a theoretical standard deviation reduction modeled as:
\begin{equation}
\text{Red}_{\sigma} = \left(1 - \sqrt{\frac{\alpha}{2-\alpha}}\right) \times 100\% \label{eq:red_sigma}
\end{equation}
